\documentclass[a4paper,11pt,fleqn,dvipsnames,twoside,openright]{memoir}  	% Openright aabner kapitler paa hoejresider (openany begge)


%%%% PACKAGES %%%%
% ¤¤ Oversaettelse og tegnsaetning ¤¤ %
\usepackage[utf8]{inputenc}					% Input-indkodning af tegnsaet (UTF8)
\usepackage[english]{babel}					% Dokumentets sprog
\usepackage[T1]{fontenc}					% Output-indkodning af tegnsaet (T1)
\usepackage{ragged2e,anyfontsize}			% Justering af elementer
%\usepackage{fixltx2e}						% Retter forskellige fejl i LaTeX-kernen

% ¤¤ Figurer og tabeller (floats) ¤¤ %
\usepackage{graphicx} 						% Haandtering af eksterne billeder (JPG, PNG, EPS, PDF)
\usepackage{multirow}                		% Fletning af raekker og kolonner (\multicolumn og \multirow)
\usepackage{multicol}         	        	% Muliggoer output i spalter
\usepackage{rotating}						% Rotation af tekst med \begin{sideways}...\end{sideways}
\usepackage{colortbl} 						% Farver i tabeller (fx \columncolor og \rowcolor)
\usepackage[table]{xcolor}					% Definer farver med \definecolor. Se mere: http://en.wikibooks.org/wiki/LaTeX/Colors
\usepackage{xcolor}

\usepackage{flafter}						% Soerger for at floats ikke optraeder i teksten foer deres reference
\let\newfloat\relax 						% Justering mellem float-pakken og memoir
\usepackage{float}							% Muliggoer eksakt placering af floats, f.eks. \begin{figure}[H]
\newcommand{\HRule}{\rule{\linewidth}{0.5mm}}
\usepackage{pdflscape}	 					% Giver mulighed for landacape mode
\usepackage{longtable}						% Giver mulighed for talber over flere side
\usepackage{subcaption}

% ¤¤ Matematik mm. ¤¤
\usepackage{amsmath}						% Avancerede matematik-udvidelser
\usepackage{amssymb}						% Avancerede matematik-udvidelser
\usepackage{stmaryrd}						% Avancerede matematik-udvidelser
\usepackage{amsthm}							% Avancerede matematik-udvidelser
\usepackage{textcomp}                 		% Symbol-udvidelser (f.eks. promille-tegn med \textperthousand )
\usepackage{mathtools}

% ¤¤ Kode indsat i rapport ¤¤ %
\usepackage{listings}						% Placer kildekode i dokumentet med \begin{lstlisting}...\end{lstlisting}




% ¤¤ Matematik funktioner ¤¤				% Sikre alle funktioner står med ordinær skriftype
	\newcommand{\e}{\text{e}}
	\newcommand{\dB}{\text{dB}}
	\newcommand{\I}{\text{I}}
	\newcommand{\V}{\text{V}}
	\newcommand{\A}{\text{A}}
	\newcommand{\W}{\text{W}}
	\newcommand{\F}{\text{F}}
	\renewcommand{\S}{\text{S}}
	\newcommand{\rms}{\text{rms}}
	\newcommand{\m}{\text{m}}
	\newcommand{\n}{\text{n}}
	\newcommand{\p}{\text{p}}
	\newcommand{\f}{\text{f}}
	\renewcommand{\k}{\text{k}}
	\newcommand{\M}{\text{M}}
	\newcommand{\G}{\text{G}}
	\newcommand{\C}{\text{C}}
	\renewcommand{\c}{\text{c}}
	\newcommand{\T}{\text{T}}
	\newcommand{\km}{\text{km}}
	\newcommand{\h}{\text{h}}
	\newcommand{\s}{\text{s}}
	\newcommand{\cm}{\text{cm}}
	\newcommand{\ms}{\text{ms}}


% ¤¤ Opsætning af enheder ¤¤ %
	\newcommand{\enhed}[1]{\hfill\left[\mathrm{ #1}\right]} 	% Skriv \enhed{"din enhed"} og dem bliver rykket til højre

% ¤¤ Referencer og kilder ¤¤ %

\usepackage[square, numbers, comma, sort&compress]{natbib}
\bibliographystyle{ieeetr}
% Udseende af litteraturlisten.

% ¤¤ Textlable setup ¤¤ %
\makeatletter
\newcommand*{\textlabel}[2]{				% Gør \textlabel mulig
  \edef\@currentlabel{#1}					% Set target label
  \phantomsection							% Correct hyper reference link
  #1\label{#2}}								% Print and store label

% ¤¤ Misc. ¤¤ %
\usepackage{lipsum}							% Dummy text \lipsum[..]
\usepackage[shortlabels]{enumitem}			% Muliggoer enkelt konfiguration af lister
\usepackage{pdfpages}						% Goer det muligt at inkludere pdf-dokumenter med kommandoen \includepdf[pages={x-y}]{fil.pdf}
\pdfoptionpdfminorversion=6					% Muliggoer inkludering af pdf dokumenter, af version 1.6 og hoejere
\pretolerance=2500 							% Justering af afstand mellem ord (hoejt tal, mindre orddeling og mere luft mellem ord)

% Kommentarer og rettelser med \fxnote. Med 'final' i stedet for 'draft' udloeser hver note en error i den faerdige rapport.
\usepackage[footnote,draft,danish,silent,nomargin]{fixme}
\renewcommand{\thefootnote}{\arabic{footnote}}

%%%% CUSTOM SETTINGS %%%%

% ¤¤ Marginer ¤¤ %
\setlrmarginsandblock{3.5cm}{2.5cm}{*}		% \setlrmarginsandblock{Indbinding}{Kant}{Ratio}
\setulmarginsandblock{2.0cm}{3.0cm}{*}		% \setulmarginsandblock{Top}{Bund}{Ratio}
\checkandfixthelayout 						% Oversaetter vaerdier til brug for andre pakker

%	¤¤ Afsnitsformatering ¤¤ %
\setlength{\parindent}{0mm}           		% Stoerrelse af indryk
\setlength{\parskip}{3mm}          			% Afstand mellem afsnit ved brug af double Enter
\linespread{1}								% Linie afstand

% ¤¤ Indholdsfortegnelse ¤¤ %
\setsecnumdepth{subsection}		 			% Dybden af nummerede overkrifter (part/chapter/section/subsection)
\maxsecnumdepth{subsection}					% Dokumentklassens graense for nummereringsdybde
\settocdepth{subsection} 						% Dybden af indholdsfortegnelsen

% ¤¤ Lister ¤¤ %
\setlist{
  topsep=0pt,								% Vertikal afstand mellem tekst og listen
  itemsep=-1ex,								% Vertikal afstand mellem items
}

% ¤¤ Visuelle referencer ¤¤ %
\usepackage[colorlinks]{hyperref}			% Danner klikbare referencer (hyperlinks) i dokumentet.
\hypersetup{colorlinks = true,				% Opsaetning af farvede hyperlinks (interne links, citeringer og URL)
    linkcolor = black,
    citecolor = black,
    urlcolor = blue
}

% ¤¤ Opsaetning af figur- og tabeltekst ¤¤ %
\captionnamefont{\small\bfseries\itshape}	% Opsaetning af tekstdelen ('Figur' eller 'Tabel')
\captiontitlefont{\small}					% Opsaetning af nummerering
\captiondelim{. }							% Seperator mellem nummerering og figurtekst
\hangcaption								% Venstrejusterer flere-liniers figurtekst under hinanden
\captionsetup{width=1\textwidth}
\setlength{\belowcaptionskip}{0pt}			% Afstand under figurteksten


% ¤¤ Opsætning til caption i listings ¤¤ %
\DeclareCaptionFormat{listing}{
  \colorbox[RGB]{255,255,255}{
    \parbox{\textwidth}{\centering#1#2#3}
  }
}
\captionsetup[lstlisting]{ format=listing, singlelinecheck=false, margin=0pt, font={normalsize} }
\renewcommand{\lstlistingname}{Code}		% Gør at der står kode og ikke listing i captionnamefont

% ¤¤ Opsaetning af listings ¤¤ %
\definecolor{commentGreen}{RGB}{34,139,24}	% Farver der bliver brugt senere
\definecolor{stringPurple}{RGB}{208,76,239}	% farver der bliver brugt senere

\lstset{language=VHDL,						% Sprog
	basicstyle=\ttfamily\small,				% Opsaetning af teksten
	keywords={for,if,while,else,elseif,		% Noegleord at fremhaeve
			  end,when,return,case,
			  switch,function,process,begin,
			  JUMP,test,ORG,in, out,
			  COMP, INT, DINT,ret
			  or, and, call, load, or,
			  then, elsif, add, sub},
	keywordstyle=\color{blue},				% Opsaetning af noegleord
	commentstyle=\color{commentGreen},		% Opsaetning af kommentarer
	stringstyle=\color{stringPurple},		% Opsaetning af strenge
	showstringspaces=false,					% Mellemrum i strenge enten vist eller blanke
	numbers=left, numberstyle=\tiny,		% Linjenumre
	extendedchars=true, 					% Tillader specielle karakterer
	columns=flexible,						% Kolonnejustering
	breaklines, breakatwhitespace=true,		% Bryd lange linjer
	captionpos=b,							% Sæt caption i bunden
	xleftmargin=14pt,						% Start margin x pt til højre(+) eller venstre(-)
}

% ¤¤ Navngivning ¤¤ %
\addto\captionsdanish{
	\renewcommand\appendixname{Appendix}
	\renewcommand\contentsname{Table of contest}
	\renewcommand\appendixpagename{Appendix}
	\renewcommand\appendixtocname{Appendix}
	\renewcommand\cftchaptername{\chaptername~}				% Skriver "Kapitel" foran kapitlerne i indholdsfortegnelsen
	\renewcommand\cftappendixname{\appendixname~}			% Skriver "Appendix" foran appendiks i indholdsfortegnelsen
}

% ¤¤ Kapiteludssende ¤¤ %
\definecolor{numbercolor}{gray}{0.7}		% Definerer en farve til brug til kapiteludseende
\newif\ifchapternonum

\makechapterstyle{jenor}{					% Definerer kapiteludseende frem til ...
  \renewcommand\beforechapskip{0pt}
  \renewcommand\printchaptername{}
  \renewcommand\printchapternum{}
  \renewcommand\printchapternonum{\chapternonumtrue}
  \renewcommand\chaptitlefont{\fontfamily{pbk}\fontseries{db}\fontshape{n}\fontsize{20}{30}\selectfont\raggedleft}
  \renewcommand\chapnumfont{\fontfamily{pbk}\fontseries{m}\fontshape{n}\fontsize{1in}{0in}\selectfont\color{numbercolor}}
  \renewcommand\printchaptertitle[1]{%
    \noindent
    \ifchapternonum
    \begin{tabularx}{\textwidth}{X}
    {\let\\\newline\chaptitlefont ##1\par}
    \end{tabularx}
    \par\vskip-2.5mm\hrule
    \else
    \begin{tabularx}{\textwidth}{Xl}
    {\parbox[b]{\linewidth}{\chaptitlefont ##1}} & \raisebox{-15pt}{\chapnumfont \thechapter}
    \end{tabularx}
    \par\vskip2mm\hrule
    \fi
  }
}											% ... her

\chapterstyle{jenor}						% Valg af kapiteludseende - Google 'memoir chapter styles' for alternativer

% ¤¤ Sidehoved ¤¤ %

\makepagestyle{AAU}							% Definerer sidehoved og sidefod udseende frem til ...
\makepsmarks{AAU}{%
	\createmark{chapter}{left}{shownumber}{}{. \ }
	\createmark{section}{right}{shownumber}{}{. \ }
	\createplainmark{toc}{both}{\contentsname}
	\createplainmark{lof}{both}{\listfigurename}
	\createplainmark{lot}{both}{\listtablename}
	\createplainmark{bib}{both}{\bibname}
	\createplainmark{index}{both}{\indexname}
	\createplainmark{glossary}{both}{\glossaryname}
}
\nouppercaseheads											% Ingen Caps oenskes


\makeevenhead{AAU}{Group 735}{}{\leftmark}					% Definerer lige siders sidehoved (\makeevenhead{Navn}{Venstre}{Center}{Hoejre})
\makeoddhead{AAU}{\rightmark}{}{Aalborg University}			% Definerer ulige siders sidehoved (\makeoddhead{Navn}{Venstre}{Center}{Hoejre})
\makeheadrule{AAU}{\textwidth}{0.5pt}						% Tilfoejer en streg under sidehovedets indhold
\makefootrule{AAU}{\textwidth}{0.5pt}{1mm}					% Tilfoejer en streg under sidefodens indhold

\copypagestyle{AAUchap}{AAU}								% Sidehoved for kapitelsider defineres som standardsider, men med blank sidehoved
\makeoddhead{AAUchap}{}{}{}
\makeevenhead{AAUchap}{}{}{}
\makeheadrule{AAUchap}{\textwidth}{0pt}
\aliaspagestyle{chapter}{AAUchap}							% Den ny style vaelges til at gaelde for chapters
															% ... her

\pagestyle{AAU}												% Valg af sidehoved og sidefod


%%%% CUSTOM COMMANDS %%%%

% ¤¤ Specielle tegn ¤¤ %
\newcommand{\decC}{^{\circ}\text{C}}
\newcommand{\dec}{^{\circ}}
\newcommand*{\rom}[1]{\expandafter\@slowromancap\romannumeral #1@}


%%%% ORDDELING %%%%

\hyphenation{}

%%%% GlOSSARIES %%%%
\usepackage[toc,nowarn]{glossaries}
\makeglossaries

%% \newacronym{label}{short name}{long name}


\newacronym{DOF}{DOF}{degrees of freedom}
\newacronym{DH}{DH}{Denavit Hartenberg}
\newacronym{JSON}{JSON}{JavaScript Objet Notation}
\newacronym{ROS}{ROS}{Robotic operating system}
\newacronym{UDP}{UDP}{User Datagram Protocol}



%% \newglossaryentry{label}{name={the real name}, description={robots are the better humans}}

\newglossaryentry{Liion}{name={LiPo batteries}, description={Lithium-Ion Polymer Battery are a rechargeable battery type which composition consist of lithium and iron. They are greatly used in the RC, automotive, industrial, electronics and alike. They are used because of their high power density compared to other batteries \cite{Lithium}.}}



%%%% COOL EXTRA FEATURES %%%%%
\usepackage{lastpage}
\usepackage{todonotes}
\newcommand{\todotom}[1]{\todo[color=red!40, author=Anders]{#1}}
\graphicspath{{./rapport/pictures/}}
\newcommand{\todohusk}[1]{\todo[color=red, author=HUSK:]{#1}}
\graphicspath{{./rapport/pictures/}}
\usepackage{epstopdf}


%%%% WARNING HACKS %%%%
\hfuzz=\maxdimen
\tolerance=100000
\hbadness=\maxdimen
\vbadness=\maxdimen
\vfuzz=\maxdimen


%%%% NICE LOOKING REFERENCE %%%%

\newcommand{\itref}[1]{\textit{\ref{#1}}}
\newcommand{\chapref}[1]{Chapter \textit{\ref{#1}: \nameref{#1}}}
\newcommand{\secref}[1]{Section \textit{\ref{#1}: \nameref{#1}}}
\newcommand{\figref}[1]{\emph{Figure \ref{#1}}}
\newcommand{\appref}[1]{\emph{Appendix: \ref{#1}}}
\newcommand{\tabref}[1]{\emph{Table: \ref{#1}}}
\newcommand{\coderef}[1]{\emph{Listings: \ref{#1}}}
\renewcommand{\eqref}[1]{\emph{Equation: (\ref{#1})}}


%%%% TIKZ MAGIC %%%%
\usepackage{schemabloc}
\usetikzlibrary{circuits}
\usepackage{tikz}
\usetikzlibrary{shapes,arrows}
\usepackage{pgfplots}
\usetikzlibrary{plotmarks}
\pgfplotsset{compat=newest}
\pgfplotsset{filter discard warning=false}
\usepackage[americanresistors,americaninductors,american voltages]{circuitikz}




%%%% CDO FIX %%%%
\usepackage{silence}
\WarningsOff[pgfplots]
\WarningFilter{latex}{Marginpar on page}
\WarningFilter{memoir}{You are using the caption package with the memoir class}
\WarningsOff[latex]
\usepackage{breakurl}
	
\raggedbottom
\begin{document}

%%% Tikz Magic %%%
%%% Tikz Magic %%%
\tikzstyle{block} = [draw, fill=blue!20, rectangle, minimum height=3em, minimum width=6em]
\tikzstyle{Integrator} = [draw, fill=blue!20, rectangle, minimum height=5mm, minimum width=5mm]
\tikzstyle{Twolineblock} = [draw, fill=blue!20, rectangle, minimum height=3em, minimum width=6em, text width = 6em, align=center]       
\tikzstyle{sum} = [draw, fill=blue!20, circle, node distance=1cm]
\tikzstyle{input} = [coordinate]
\tikzstyle{output} = [coordinate]
\tikzstyle{pinstyle} = [pin edge={to-,thin,black}]
\tikzstyle{box} = [draw,rounded corners, minimum height=15mm, minimum width=20mm, align=center, text centered]
\tikzstyle{BlackBox} = [draw, fill=black, rounded corners, minimum height=15mm, minimum width=20mm, align=center, text=white, text centered]
\tikzstyle{FlowIF} = [diamond, draw, fill=blue!20, text width=8.5em, text badly centered, node distance=3cm, inner sep=0pt,align=center,aspect=3]
\tikzstyle{FlowBlock} = [rectangle, draw, fill=blue!20, text width=8em, text centered, rounded corners, minimum height=3em]
\tikzstyle{FlowCloud} = [draw, ellipse,fill=black!40!red,box, node distance=3cm, minimum height=2em]
\tikzstyle{TestBox} = [rectangle,draw, fill=black!20, minimum height=8mm, minimum width=8mm]
\tikzstyle{TestDiamond} = [diamond, draw, fill=black!20, aspect=2, minimum height=8mm]
%\tikzstyle{TestDiamond} = [diamond, draw, fill=black!20, minimum height=5mm, minimum width=5mm]

\tikzstyle{TestBox} = [rectangle,draw, fill=black!20, minimum height=8mm, minimum width=8mm]
\tikzstyle{TestCircle} = [circle, draw, fill=black!20, minimum height=6mm, minimum width=6mm]
\tikzstyle{TestTable} = [rectangle, draw, fill=black!60, minimum height=0.01mm, minimum width=9mm, text=white]
\tikzstyle{TestBoxSmall} = [rectangle,draw, fill=black!, minimum height=8mm, minimum width=8mm,text=white]
\tikzstyle{LegendBox} = [rectangle,draw, minimum height=7mm, minimum width=20mm]
\tikzstyle{Sysbox} = [draw,rounded corners, minimum height=15mm, minimum width=9em, align=center, text centered,text width = 10.5em]
\tikzstyle{SysBlackBox} = [draw, fill=black, text=white, rounded corners, minimum height=15mm, minimum width=9em, align=center, text centered,text width = 10em]


% Forindhold - Kun romer tal

\frontmatter

%Formalia
%%main for introduktion

%\includepdf[pages={1}]{rapport/introduktion/forsiden.pdf}

%\includepdf[pages={1}]{forside/forsiden.pdf}

%\thispagestyle{empty}

\begin{center}

\vspace*{\fill}

\textsc{\LARGE Aalborg University}\\[1.0cm]

\HRule \\[0.4cm]
{ \HUGE \bfseries  Teleoperation of surgical robot using force feedback \\[0.5cm] }

\HRule \\[1.5cm]%


\begin{minipage}{0.4\textwidth}
\begin{flushleft} \large
Electronic \& IT:\\
Control and Automation:\\
Fourth year of study
\end{flushleft}
\end{minipage}
\begin{minipage}{0.4\textwidth}
\begin{flushright} \large
Group: \\
CA-735
\end{flushright}
\end{minipage}

\vspace*{\fill}

\textsc{\Large Worksheet}\\[1.0cm]

{\large 20. December 2016}

\end{center}
\newpage

%\phantomsection
\pdfbookmark[0]{Titelblad}{titelblad}
\thispagestyle{empty}

\begin{minipage}[t]{0.48\textwidth}
\vspace*{-25pt}			%\vspace*{-9pt}

\begin{figure}[H] 
\includegraphics[width=0.95\textwidth]{rapport/introduktion/aau_logo1}
\end{figure} 
\end{minipage}
\hfill
\begin{minipage}[t]{0.48\textwidth}
{\small 
\textbf{Third year of study}  \\
Electronic og IT \\
Fredrik Bajers Vej 7 \\
DK-9220 Aalborg East, Denmark\\
http://www.es.aau.dk}
\end{minipage}


\vspace*{1cm}

\begin{minipage}[t]{0.48\textwidth}
\textbf{Topic:} \\[5pt]\bigskip\hspace{2ex}BSc Project Control Engineering 

\textbf{Project:} \\[5pt]\bigskip\hspace{2ex}
P6-project

\textbf{Project time:} \\[5pt]\bigskip\hspace{2ex}
February 2016 - May 2016

\textbf{Projectgroup:} \\[5pt]\bigskip\hspace{2ex}
16gr633	

\textbf{Participants:} \\[5pt]\hspace*{2ex}
Daniel Bähner Andersen \\\hspace*{2ex}
Nicolaj Vinkel Christensen \\\hspace*{2ex}
Ralf Victor Lømand R. Christiansen\\\hspace*{2ex}
Simon Bjerre Krogh \\\hspace*{2ex}
Thomas Holm Pilgaard \\

\textbf{Supervisor:} \\[5pt]\hspace*{2ex}
Tom S. Pedersen \\\hspace*{2ex}
Kirsten M. Nielsen \\\bigskip\hspace{2ex}

\vspace*{3.5cm}

\textbf{Circulation: 8} \\
\textbf{Number of pages: 152}\\
\textbf{Appendix: 73 + CD} \\
\textbf{Completed 25-05-2016}\\
\end{minipage}
\hfill
\begin{minipage}[t]{0.483\textwidth}
\textbf{Synopsis:} \\[5pt]
\fbox{\parbox{7cm}{\bigskipThis project covers the development of a controller system for a model container crane. The goal of this project is to automatically relocate containers fast, reliably and efficiently. To achieve this, the hardware platform has been improved and is now based on digital motor drivers together with a FPGA as the main controlling unit. The control system is based on multiple controllers derived by the use of classical control methods, namely the root locus method. The system models describing the behavior of the crane are derived using a combination of Newtons law's and free body diagrams to describe this graphically. An issue discovered during the project, is that the static friction is not constant along the axes of the setup. In order to compensate for the static friction a current offset is implemented. However, this yields issues as the static friction is not constant. The controller system for relocating containers has been proved functional.Additionally an interface has been implemented, this allows to relocate containers from eight different positions along the x-axis and three locations along the y-axis.          \bigskip}}
\end{minipage}

\vfill

%{\footnotesize\itshape Rapportens indhold er frit tilgængeligt, men offentliggørelse (med kildeangivelse) må kun ske efter aftale med forfatterne. Forside billed }

\chapter*{Preface}

\textcolor{red}{This project comprises of implementing a functional controller system for the provided miniature model crane. This is conducted with the purpose of relocating containers. A basic understanding of electronics is expected of the reader. The report is constructed with different chapters, each covering a main part of the entire project. First, a thorough analysis of the crane setup and the modernizations carried out to the hardware platform is presented. Furthermore, the construction of models describing the crane and verification of these is presented. The second part consists of designing controllers based on requirements derived from an earlier chapter. The controllers are implemented and tested. Finally an acceptance test is performed and this is documented in the appendix. All of the raw test information, test programs and relevant datasheets are stored on the enclosed CD. Regarding figures without a source reference, these are created specifically for this report. A big thanks to the group's supervisors Tom S. Pedersen and Kirsten M. Nielsen for help with discussing possible solutions and giving insight into the current solutions in use, concerning classical control theory. Thanks to Rasmus Viking Lømand Ravgård Christiansen, bachelor in mechanical engineering, for helping creating a \textbf{Solid Work} \cite{Solid_Work} model for parts of the crane. }

\begin{flushright}
Aalborg University, 20th of December 2016
\end{flushright}





\vfill

\begin{table}[H]
	\centering
		\begin{tabular}{c c }
			\underline{\phantom{mmmmmmmmmmmmmmmmmmm}}       & \underline{\phantom{mmmmmmmmmmmmmmmmmmm}} \\
			Dániel Bolgár			 & Simon Bjerre Krogh  \\
			\textit{dbolga16,@student.aau.dk} & \textit{skrogh13,@student.aau.dk}\\
			&\\
			&\\
			\underline{\phantom{mmmmmmmmmmmmmmmmmmm}}       & \underline{\phantom{mmmmmmmmmmmmmmmmmmm}} \\
		    Filip Marić			 & Nicolas Silviani\\
			\textit{fmaric16,@student.aau.dk} & \textit{nsilva16@student.aau.dk} \\
			&\\
			&\\	
		\end{tabular}
\end{table}




 %%%%%%%     <->    Front page, synopsis and preface included in this
\section*{The four Ws'} % \label{app:...}
Please make a new document for each topic, so it will be easy to edit if there is a mistake.

\begin{itemize}
\item Who are you addressing: 

scientists who are specialists in your field
of research, a wider group of scientists, fellow students, or public
audiences?
\item Why is your message important? 

Why are your communicating it?
Presumably you are not doing it just for credits, but to add to the
pool of knowledge
\item What are your main findings or “take-home” messages? '

What are
you going to present - new research results or a review of a topic?
What prior knowledge, expectations and questions might your
audience have? What technical language do they understand?
\item How can you best deliver your message and satisfy the audience’s need? 

How will the audience use its new knowledge?

\end{itemize}
\section*{Test template} % \label{app:...}
Use this template when you need to make a test of somthing

\subsection*{Purpose:}

\subsection*{Test equipment:}

\subsection*{Procedure:}

\subsection*{Measuring data:}

\subsection*{Results:}

\subsection*{Uncertainties of measurement:}

\subsection*{Conclusion:}


\tableofcontents*

% Hovedindhold - nummereres fra ''side 1 af sidste side''b
\mainmatter 
\makeevenfoot{AAU}{\thepage \text{} of \pageref{LastPage}}{}{}
\makeoddfoot{AAU}{}{}{\thepage \text{} of \pageref{LastPage}}
%%%%%%%%%%%%%%%%%%%%%%%%%%%%%%%%%%%%%%%%%%%%%%%%%%%%%%%%%%%%%%%
% Delete this before turning in


%%%%%%%%%%%%%%%%%%%%%%%%%%%%%%%%%%%%%%%%%%%%%%%%%%%%%%%%%%%%%%%
%  Include sections here   %

\chapter{Kinematic for the Geomagic touch} % \label{app:...}

bla bla bla

\begin{table}[h!]
\centering
\label{my-label}
\begin{tabular}{|l|l|l|l|l|}
\hline
 &  &  &  &  \\ \hline
 &  &  &  &  \\ \hline
 &  &  &  &  \\ \hline
 &  &  &  &  \\ \hline
 &  &  &  &  \\ \hline
 &  &  &  &  \\ \hline
 &  &  &  &  \\ \hline
 &  &  &  &  \\ \hline
 &  &  &  &  \\ \hline
\end{tabular}
\caption{Pleas give me some kind of information!}
\end{table}



All surgical interventions with the Da Vinci robot require some sort of feedback for the user, in our work we are covering force/haptic feedback.
This type of feedback translates the forces acting on the surgical instrument to the haptic controller in some manner.
The haptic controller used in this work is the Phatnom Omni (aka Geomagic Touch) developed by Sensable Technologies.
It is one of the most cost effective haptic controllers currently on the market.

In order for useful feedback to be created, we need to be able to control the direction of the force created by the Phantom Omni.
A logical first step would be to derive the relationship between joint coordinates and end-effector position (the end-effector in this case being the point of contact between the user and device). 
Once we have this, it is possible to develop other constructs such as Jacobian and dynamic equations.

We derive the forward kinematics of the Phantom omni using the DH-algorithm for simplicity. Also, since the actuated joints only control the position, and the orientation joints are concentrated in a small space, we can separate the FK matrices into translational and rotational parts.

\begin{align}
\begin{tabular}{ |p{2cm}|p{2cm}|p{2cm}|p{2cm}|  }
 \hline
 \multicolumn{4}{|c|}{Phantom Omni DH parameters} \\
 \hline
 a &d &$\theta$ &$\alpha$ \\
 \hline
 0          & 0   &$q_{1}$&   0\\
 0          &   0           & $q_{2}$ &$\frac{\pi}{2}$\\
 A         & 0   & $q_{3}$    &  0\\
 0          & D             & $q_{4}$     & $\frac{\pi}{2}$\\
 0& 0 & $q_{5}$ &-$\frac{\pi}{2}$\\
  0& 0 & $q_{6}$ &$\frac{\pi}{2}$\\
 \hline
\end{tabular}
\end{align}



\begin{equation}
\mathbf{T_0^1} =\begin{bmatrix} \cos\theta_{1}& 0 & \sin\theta_{1} & 0\\
\sin\theta_{1}& 0 & -\cos\theta_{1} & 0\\
0 & 1 & 0 & d_{1}\\
0 & 0 & 0 & 1\\
\end{bmatrix}
\end{equation}







%%%%%%%%%%%%%%%%%%%%%%%%%%%%%%%%%%%%%%%%%%%%%%%%%%%%%%%%%%%%%%%
%  Bilag  %
\appendix

\clearpage

\bibliography{bibliography}

\end{document}